\section{Introduction}
\label{sec:intro}

Two decisions share the same scarce vehicle time on a high-frequency bus corridor. The dispatch planner can move scheduled departures toward periods of higher demand without adding daily trips. At a stop, holding can improve spacing, but it delays riders already on board and can postpone a vehicle's return for a trip in the opposite direction. A stop-level action that looks helpful on a headway metric need not be helpful to passengers or to the next terminal departure \citep{osuna1972control,daganzo2009headway,berrebi2018comparing}.

We study a fixed fleet circulating between two terminals. FreqDuet uses an upper policy to adjust a departure plan and a trip budget retention parameter every 15 minutes. A lower policy acts after a stop's boarding exchange. It can hold only within the remaining trip budget. The budget is based on a projection of which vehicles can cover upcoming departures, so it constrains actions without assuming that forecast return times will be realized. The upper and lower policies also use different summaries of a shared APC-derived signal: the upper policy emphasizes smooth directional demand and the lower policy emphasizes recent local departures from that trend.

The study makes three contributions:
\begin{enumerate}
  \item It specifies an executable joint dispatch and holding architecture under a fixed fleet and trip roster. An analytical two-gap model explains why unequal planned spacings, onboard load and future vehicle availability matter to the holding decision.
  \item It implements quota-aware dispatch and per-trip holding budgets in a discrete, two-level learning controller. The experiments compare the whole policy with fixed dispatch, two named headway feedback rules, a rolling dispatch heuristic, and other learned policies.
  \item It separates two empirical questions: whether joint control improves passenger outcomes relative to fixed dispatch, and whether the tested demand-feature representation improves a shared two-level controller. We evaluate paired simulated demand days and report training-seed sensitivity and the operating regimes where the comparison changes.
\end{enumerate}

The two nominal-corridor tests show shorter passenger journeys relative to no holding, mainly from less waiting, with some loss of arrival regularity and longer readiness delay. Comparisons with dynamic alternatives are more mixed. The MBTA-informed simulation supports a benefit from the tested demand features relative to a coarse APC feature block. The synthetic tidal corridor does not support the same conclusion. These distinctions matter when choosing between a joint learned controller and a simpler operating rule.
